\scanpage{349}
% top centre
\mspage{265}
% lines 1–4
\uline{Corollaire} 4.10.9 : Soit $N\unsure{\mathrel{-}|\longrightarrow}J$ une ``cofibration en paratopos'',$^*$ i.e. un foncteur cofibrant tel que les fibres soient des cat. paratopos, et les foncteurs
% upper left beside opening quotation
% slanted
\margin{\unsure{terme|termes} \unsure{ambigu|ambigus}, cf la double \unsure{explicitation}}
% lines 5–10
\unsure{cochgt|cochgt.} de base \deleted{\unsure{soient}}\inserted{soient des foncteurs ``image \deleted{\unsure{direc}}\unsure{réciproque|inverse}'' \unsure{ou|,} (\unsure{resp.} ``image \unsure{directe|inverse}''),} pour des morphismes de paratopos, i.e qu'ils commutent aux (petites) $\varinjlim$\inserted{(\unsure{resp.|ainsi} aux petites $\varprojlim$)}. Supposons de plus $J$ \unsure{équivalent|équivalente} à une petite catégorie. Alors $\underline{\unsure{\Gamma|\mathscr T}}(N/J)$ est une \inserted{catégorie} paratopos.

% lines 11–17
En effet, c'est une catégorie accessible en vertu de 4.10.8, d'autre part l'hypothèse sur les $\lambda_*$ implique aussitôt que \deleted{\unsure{la}} $\underline{\Gamma}(N/J)$ est stable par petites $\varprojlim$ (\unsure{resp} par petites $\varprojlim$). Donc $\underline{\Gamma}(N/J)$ est une catégorie paratopos en vertu de 4.\unsure{5|8}.9 $b\Longrightarrow a$ (resp. $c\Rightarrow\unsure{o|a|0}$).

% lines 18–20
\uline{Corollaire} 4.10.10. Soit $N\mathrel{-}J$ une catégorie \uline{fibrée} (non cofibrée !), telle que les fibres soient des catégories paratopos et les foncteurs
% archive sheet number: 348
